\documentclass[11pt, oneside]{article}   	% use "amsart" instead of "article" for AMSLaTeX format
\usepackage[margin=1 in]{geometry}                		% See geometry.pdf to learn the layout options. There are lots.
\geometry{letterpaper}                   		% ... or a4paper or a5paper or ... 
%\geometry{landscape}                		% Activate for rotated page geometry
%\usepackage[parfill]{parskip}    		% Activate to begin paragraphs with an empty line rather than an indent
\usepackage{graphicx}				% Use pdf, png, jpg, or eps§ with pdflatex; use eps in DVI mode
								% TeX will automatically convert eps --> pdf in pdflatex		
\usepackage{amssymb}

\usepackage{amsmath}

\usepackage{url}
\usepackage{hyperref}

\newtheorem{proposition}{Proposition}
\newtheorem{example}{Example}

\usepackage{listings}
%SetFonts

%SetFonts
\usepackage{xcolor}
\usepackage{float}

\title{\vspace{-.3 in} Operations in Action}
\author{Nick Arnosti}
\date{August 6, 2018}							% Activate to display a given date or no date

\begin{document}
\maketitle


Many people are familiar with the fact that legroom in airplane seats has been declining over the years. Some people are upset about this fact. A few even sued the Federal Aviation Administration, arguing that the cramped quarters increase the risk during an emergency evacuation. Setting aside that I generally believe that emergency safety measures on planes are not cost-effective (emergencies are rare, and emergencies in which life vests, safety slides, etc are actually helpful appear to be almost non-existent), the argument sounds logical -- doesn't it take longer to get in and out of a cramped seat than a spacious one?

Turns out, that's not the right question to ask. Instead, ask ``What is the bottleneck?" Evacuation time is dominated by the time required to get every passenger through the exits. Make it faster to stand up, and passengers will simply spend more time waiting after standing. On this basis, the FAA {\color{blue} \underline{\href{https://www.npr.org/2018/07/05/626090518/faa-to-scrunched-passengers-sardine-seats-won-t-be-regulated}{won the case}}}.

One final thought. The whole idea of less legroom is to pack more people into one plane, which could in fact increase evacuation time. But then the regulation should specify a maximum number of passengers per exit door, not a minimum seat pitch.


\end{document}